\chapter{WF-009 - Keep Release PRs Under Automated Review File Limits}\label{record-205}

\noindent{\footnotesize\color{BookMuted}\path{docs/design/WF-009-keep-release-prs-under-automated-review-file-limits.md}\par}\smallskip

Legend: \href{https://github.com/flyingrobots/bijou/blob/1853ad902f6bf219326409e7f743a5a1d6517dbc/docs/legends/WF-workflow-and-delivery.md}{WF - Workflow and Delivery}

\section{Sponsor human}\label{record-205-sponsor-human}

A maintainer preparing release-boundary work that may span many packages, docs, examples, tests, and backlog notes.

\section{Sponsor agent}\label{record-205-sponsor-agent}

An agent executing the METHOD loop through pull, implementation, commit, push, PR, and merge without producing a PR too large for automated review.

\section{Hill}\label{record-205-hill}

Release-boundary work is sliced into reviewer-sized PRs before the final release sync, so automated review, CI, and human inspection all see bounded changes.

\section{Why this cycle exists}\label{record-205-why-this-cycle-exists}

CodeRabbit skipped PR \#65 because the diff exceeded its 150-file cap. PR \#69 also reached the danger zone: the local pre-push hook warned about 23 commits, and CodeRabbit selected 149 files for processing.

The issue is not that large releases are bad. The issue is that a large release must be composed from smaller reviewable precursor PRs instead of arriving as one oversized branch.

\section{Playback questions}\label{record-205-playback-questions}

\begin{enumerate}
\def\labelenumi{\arabic{enumi}.}
\tightlist
\item
  Can an operator identify when a release branch must be split before PR?
\item
  Can an agent choose the next release-boundary slice without relying on chat memory?
\item
  Can a reviewer see whether a PR is intentionally a precursor slice or a final release sync?
\item
  Can the final release PR stay mostly about version, changelog, long-form release docs, and release validation rather than hiding unrelated features?
\end{enumerate}

\section{Accessibility / assistive reading posture}\label{record-205-accessibility-assistive-reading-posture}

This is workflow documentation. It must remain useful in plain markdown, terminal readers, MCP document summaries, and copied text. The policy uses short numbered rules instead of relying on diagrams.

\section{Localization / directionality posture}\label{record-205-localization-directionality-posture}

No localized runtime behavior changes. The guidance should avoid idioms that make the release process harder to translate later.

\section{Agent inspectability / explainability posture}\label{record-205-agent-inspectability-explainability-posture}

The slicing threshold and required response live in repo docs, not chat memory:

\begin{itemize}
\tightlist
\item
  \href{https://github.com/flyingrobots/bijou/blob/cbd70612140e7a450d764b3c7be1c6af878b99bc/docs/WORKFLOW.md}{docs/WORKFLOW.md}
\item
  \href{https://github.com/flyingrobots/bijou/blob/cbd70612140e7a450d764b3c7be1c6af878b99bc/docs/release.md}{docs/release.md}
\item
  this cycle doc
\end{itemize}

\section{Non-goals}\label{record-205-non-goals}

\begin{itemize}
\tightlist
\item
  changing branch protection
\item
  disabling CodeRabbit
\item
  forcing every tiny fix into its own PR
\item
  tagging or publishing \texttt{v\allowbreak{}6\allowbreak{}.\allowbreak{}0\allowbreak{}.\allowbreak{}0}
\end{itemize}

\section{Decision}\label{record-205-decision}

Bijou release work should use precursor PRs whenever a branch is likely to exceed automated review limits.

The practical thresholds are:

\begin{itemize}
\tightlist
\item
  hard stop: more than 140 changed files before opening a PR
\item
  caution: more than 10 commits on one branch
\item
  caution: unrelated legends or cycle families mixed in one branch
\item
  caution: generated snapshots or docs mixed with unrelated runtime changes
\end{itemize}

When a threshold is hit, split by cycle or product surface:

\begin{itemize}
\tightlist
\item
  runtime engine
\item
  TUI shell and interaction
\item
  core components
\item
  DOGFOOD and docs proof
\item
  backlog/release-lane bookkeeping
\item
  release version/changelog/tag prep
\end{itemize}

The final release PR should be boring. It should primarily contain version bump, changelog, release docs, release-readiness proof, and any final metadata adjustments that cannot land earlier.

\section{Implementation outline}\label{record-205-implementation-outline}

\begin{enumerate}
\def\labelenumi{\arabic{enumi}.}
\tightlist
\item
  Add PR sizing discipline to the short workflow guide.
\item
  Add release-boundary slicing guidance to the release process.
\item
  Update stale release facts while editing the release process so the docs point at \texttt{v\allowbreak{}6\allowbreak{}.\allowbreak{}0\allowbreak{}.\allowbreak{}0} as the next shaping lane.
\item
  Keep the v6 backlog lane linked to this pulled cycle.
\end{enumerate}

\section{Tests / validation}\label{record-205-tests-validation}

\begin{itemize}
\tightlist
\item
  \texttt{n\allowbreak{}p\allowbreak{}m\allowbreak{} \allowbreak{}r\allowbreak{}u\allowbreak{}n\allowbreak{} \allowbreak{}d\allowbreak{}o\allowbreak{}c\allowbreak{}s\allowbreak{}:\allowbreak{}i\allowbreak{}n\allowbreak{}v\allowbreak{}e\allowbreak{}n\allowbreak{}t\allowbreak{}o\allowbreak{}r\allowbreak{}y}
\item
  \texttt{g\allowbreak{}i\allowbreak{}t\allowbreak{} \allowbreak{}d\allowbreak{}i\allowbreak{}f\allowbreak{}f\allowbreak{} \allowbreak{}-\allowbreak{}-\allowbreak{}c\allowbreak{}h\allowbreak{}e\allowbreak{}c\allowbreak{}k}
\item
  \texttt{n\allowbreak{}p\allowbreak{}m\allowbreak{} \allowbreak{}r\allowbreak{}u\allowbreak{}n\allowbreak{} \allowbreak{}l\allowbreak{}i\allowbreak{}n\allowbreak{}t}
\end{itemize}

\section{Retrospective}\label{record-205-retrospective}

Closed by documenting PR slicing thresholds in the workflow guide, adding release-boundary sizing rules to the release process, and updating stale release posture from \texttt{v\allowbreak{}5\allowbreak{}.\allowbreak{}0\allowbreak{}.\allowbreak{}0} to \texttt{v\allowbreak{}6\allowbreak{}.\allowbreak{}0\allowbreak{}.\allowbreak{}0}.

